\begin{textsample}{POS Dim 1 – vem\_ed\_35 – Score 90.00 – vem\_ed\_35/t187.txt}  \label{ex:f1_pos_011}
5º IVES 2026: obstáculos e \textbf{pontes} Com o tema “Cooperation in COIL Virtual Exchange: hurdles and bridges” (Cooperação em Intercâmbios Virtuais do tipo COIL: obstáculos e \textbf{pontes}), a \textbf{quinta} edição do International Virtual Exchange Symposium (IVES CGESG 2026) reuniu quatro palestrantes internacionais: Sara Pittarello (UNICollaboration, rede europeia de \textbf{apoio} aos Intercâmbios Virtuais), Carlos Maza (UNAM, México), Henry Shepherd (IVEC, EUA) e Salome Escalona (Bukidnon State University, Filipinas).
O evento online ocorreu em 28 de maio de 2026 e engajou 189 participantes.
Sara Pittarello, diretora da UNICollaboration, ressaltou \textbf{que} os Intercâmbios Virtuais \textbf{permitem} \textbf{expandir} acesso a experiências internacionais, desenvolver \textbf{competências} globais e \textbf{habilidades} interculturais, apoiam equidade e inclusão e preparam os estudantes para \textbf{ambientes} de trabalho globalizados.
Não se trata apenas de conectividade, mas de \textbf{interações} \textbf{significativas}, objetivos \textbf{compartilhados}, colaboração estruturada, aprendizagem cocriada e parcerias de \textbf{longo} prazo.
Apresentou projetos financiados pelo Erasmus+, principal \textbf{programa} de mobilidade acadêmica europeia, \textbf{que} abordam temas como “Critical Virtual Exchange in AI”, voltado para uma \textbf{pedagogia} \textbf{crítica}, a consciência ética e o letramento em Inteligência Artificial por meio de Intercâmbios Virtuais entre Europa e África subsaariana; \textbf{desenvolvimento} \textbf{profissional} de professores com tecnologias digitais e o \textbf{programa} VAMOS, voltado a Intercâmbios Virtuais \textbf{que} solucionam problemas de \textbf{desenvolvimento} \textbf{sustentável}.
O VAMOS integrou universidades da Europa, do Brasil e de Honduras.
Para encerrar, Pittarello \textbf{destacou} quatro \textbf{pontos}: - Intercâmbios Virtuais \textbf{são} mais \textbf{que} tecnologia - Sustentabilidade dos projetos requer \textbf{estrutura} e \textbf{estratégia} - Facilitação \textbf{é} fundamental, não opcional - As parcerias devem \textbf{ser} de propriedade conjunta Carlos Maza, da UNAM, \textbf{destacou} a \textbf{construção} de “\textbf{pontes} com alicerces \textbf{humanos}” por meio da abordagem COIL (Collaborative Online International Learning).
Lembrou como a pandemia de covid-19 impulsionou rapidamente a comunidade universitária para o ensino remoto, \textbf{criando} as bases para a colaboração internacional virtual.
Com a retomada das atividades presenciais, a UNAM aproveitou essa infraestrutura para \textbf{implementar} o \textbf{programa} COIL, \textbf{conectando} salas de aula \textbf{locais} a instituições ao redor do mundo.
Nesse \textbf{contexto}, a adesão a redes como COIL Connect e Red LatAM COIL \textbf{foi} \textbf{essencial} para estruturar parcerias globais.
A experiência demonstrou \textbf{que} o principal pilar \textbf{são} os docentes.
Para sua formação continuada, a universidade oferece oficinas de capacitação e eventos acadêmicos como as Jornadas COIL UNAM.
Além de formar professores, \textbf{é} importante preparar os alunos para as experiências internacionais virtuais, \textbf{reduzindo} ansiedade e \textbf{promovendo} engajamento consciente.
Maza \textbf{destaca} \textbf{que} o COIL democratiza o acesso à internacionalização, \textbf{ampliando} oportunidades em uma universidade grande como a UNAM (\textbf{que} tem mais de 370 mil alunos).
Entre os desafios, \textbf{destaca} a difusão do conceito COIL para \textbf{toda} a comunidade acadêmica e a consolidação de diretrizes \textbf{institucionais}.
Henry Shepherd, do International Virtual Exchange Consortium (IVEC), \textbf{destacou} dois obstáculos aos projetos COIL e duas \textbf{pontes} para superá-los.
O primeiro \textbf{é} a dificuldade de \textbf{reconhecer} \textbf{diferenças} de \textbf{perspectiva} entre parceiros. “Até mesmo a palavra COIL não significa a mesma \textbf{coisa} em \textbf{todos} os lugares”, \textbf{observou}.
Nossas experiências e preferências moldam nossas suposições sobre \textbf{cada} aspecto do COIL”.
Deixar de considerar o \textbf{ponto} de vista do outro \textbf{é} fonte de muitos desafios nos projetos.
Para superar esse obstáculo, Shepherd sugere colocar-se no lugar do parceiro, tomar decisões conjuntamente desde o planejamento inicial e, quando não houver meio-termo entre os parceiros, revezar: “sua plataforma, meu app; \textbf{horário} da sua aula primeiro, depois o meu”.
O segundo obstáculo \textbf{é} a incerteza.
Geralmente, apontamos os desafios óbvios: \textbf{diferenças} de fusos \textbf{horários}, de tecnologias, calendários ocupados.
O \textbf{verdadeiro} problema pode estar em outro \textbf{ponto}: \textbf{diferenças} nos estilos de comunicação, de gestão de projetos, ou de abordar uma \textbf{questão} \textbf{que} não \textbf{foi} tratada. “A comunicação \textbf{que} para uns parece \textbf{profissional}, para outros pode \textbf{ser} sentida como fria. \textbf{É} difícil ler comportamentos online em alguém com quem você nunca encontrou pessoalmente”, ponderou.
Além disso, as culturas expressam o atrito de \textbf{formas} diferentes. “Podemos simplesmente tentar superar isso e \textbf{seguir} em frente, podemos nos afastar e nos distanciar do nosso parceiro, ou ir embora sem dar explicações e deixar o projeto ir por água abaixo.
Infelizmente, tudo isso acontece no COIL”.
Por isso, Shepherd encoraja os parceiros a conversarem sobre \textbf{expectativas} sobre como vão trabalhar juntos e encarar os desafios, antes mesmo \textbf{que} aconteçam.
Reconhecer \textbf{que} as \textbf{pessoas} colaboram de \textbf{formas} diferentes: silêncio não significa desinteresse.
E resolver atritos antes \textbf{que} se prolonguem por semanas.
Outra recomendação \textbf{é} incluir nas reuniões de planejamento com o professor parceiro e na preparação para o projeto com os alunos, uma discussão sobre a complexidade e as oportunidades de aprendizagem \textbf{que} surgem ao lidar com a cooperação e a colaboração a distância.
Para \textbf{que} os docentes possam \textbf{aprender} com os pares, \textbf{fortalecendo} \textbf{práticas}, \textbf{compartilhando} experiências e soluções, Shepherd recomenda o engajamento em redes internacionais como COIL Connect, Red LatAM COIL, BraVE, UNICollaboration, AfriCOIL, Japan COIL Association. “Quando usamos a paciência e a curiosidade para compreender melhor os \textbf{pontos} de vista dos outros, construímos relacionamentos, desenvolvemos \textbf{habilidades} críticas úteis e encontramos satisfação por ter feito algo difícil, e isso realmente vale a pena fazer”.
Salome Escalona, da Bukidnon State University (BukSU, Filipinas), apresentou obstáculos e \textbf{pontes} sob o \textbf{ponto} de vista de coordenadores e professores.
Citou \textbf{que}, ao ingressar em \textbf{ambientes} globais sofisticados, repletos de \textbf{profissionais} altamente experientes, professores podem experimentar a “síndrome do impostor”.
Para superar essa síndrome, \textbf{é} importante valorizar a mentoria e a \textbf{construção} de relacionamentos, \textbf{aprendendo} com a generosidade de parceiros experientes.
Alguns professores podem se assustar com a carga de trabalho \textbf{gerada} pelos projetos COIL, \textbf{que} podem \textbf{ser} vistos como “\textbf{tarefa} extra”.
Para evitar essa visão, coordenadores de projetos COIL devem demonstrar como as atividades podem \textbf{ser} integradas aos objetivos de aprendizagem da \textbf{disciplina}.
Assim, os professores integram o projeto COIL a \textbf{tarefas} \textbf{que} já fazem \textbf{parte} do plano de ensino.
Como lidar com a ansiedade dos estudantes e barreiras de comunicação?
A professora sugere a \textbf{construção} de um \textbf{ambiente} de comunicação seguro e encorajador.
Para contornar a redução do engajamento estudantil e a fadiga digital, Salome Escalona sugere um ritmo de atividades mais espaçado e flexível, com \textbf{tarefas} ancoradas em \textbf{conexões} \textbf{humanas} reais. “Quando confrontados com a \textbf{tarefa} de \textbf{produzir} um trabalho verdadeiramente intercultural, os alunos deixaram de encarar a atividade como apenas mais uma \textbf{tarefa}”.
E conclui com uma \textbf{reflexão}: “O \textbf{intercâmbio} virtual \textbf{é} muito mais do \textbf{que} uma exigência administrativa ou uma lista de parceiros internacionais. \textbf{É} um forte lembrete de \textbf{que} as \textbf{diferenças} geográficas, culturais e linguísticas desaparecem quando priorizamos a \textbf{conexão} \textbf{humana}”.
Os quadros a \textbf{seguir} sintetizam os obstáculos e \textbf{pontes} sob as \textbf{perspectivas} dos coordenadores e professores, apresentadas por Salome Escalona (leia sobre o PCI realizado entre Fatecs Pompeia, Guarulhos e Zona Leste e BukSU no “Artigo de Opinião”).
Histórico do IVES As quatro edições \textbf{anteriores} do International Virtual Exchange Symposium (IVES), realizadas entre 2022 e 2025, \textbf{foram} registradas nos \textbf{números} 15, 18, 24 e 29 de VEm e \textbf{publicadas} em playlists do Canal CGESG no YouTube: IVES 2025: https://bit.ly/3GWupxd IVES 2024: http://bit.ly/3SQtvFi IVES 2023: https://bit.ly/4k3ymPq IVES 2022: https://bit.ly/3FnE5R9 Ao \textbf{longo} dos últimos anos, a trajetória do IVES vem \textbf{consolidando} os Intercâmbios Virtuais como \textbf{estratégia} de Internacionalização em Casa para o Ensino Superior no Centro Paula Souza, bem como a evolução dos Projetos Colaborativos Internacionais (PCIs) realizados nas Fatecs sob a \textbf{coordenação} da equipe de Apoio à Internacionalização do Ensino Superior da Coordenadoria Geral de Ensino Superior de Graduação do CPS.
A primeira edição (IVES 2022) teve caráter inaugural, reunindo parceiros internacionais (China, Índia, África do Sul e Colômbia) e apresentando PCIs desenvolvidos nas Fatecs Jaboticabal, Garça, Franco da Rocha, São José do Rio Preto, Matão e Itu.
Os relatos de experiência apresentaram a abordagem COIL (Collaborative Online International Learning) tropicalizada nos PCIs.
Também \textbf{foi} \textbf{mencionada} a criação do ProCin nas Etecs, \textbf{programa} \textbf{iniciado} em 2021 pela Coordenadoria Geral de Ensino Médio e Técnico e pela ARInter, com a mentoria da equipe dos PCIs.
O evento \textbf{destacou} o Intercâmbio Virtual como ferramenta de Internacionalização em Casa, \textbf{baseada} na colaboração remota e no \textbf{desenvolvimento} de \textbf{competências} interculturais.
Na segunda edição (IVES 2023), houve a participação de cinco especialistas internacionais e ampliação do público, \textbf{que} subiu de 176 para 190 participantes.
As palestras \textbf{foram} oferecidas por coordenadores de Intercâmbios Virtuais da DePaul University (EUA), DUOC UC (Chile), Universidad Católica Andrés Bello (Venezuela), University of Minnesota System e SUNY COIL Center (EUA).
As discussões enfatizaram tendências e desafios do COIL.
Hope Windle (SUNY COIL Center) trouxe conceitos como inclusão, interseccionalidade e “esperança \textbf{crítica}”, além de projeções para o \textbf{futuro} dos \textbf{intercâmbios} virtuais.
O evento se posiciona como espaço de \textbf{troca} global de boas \textbf{práticas}.
A \textbf{terceira} edição (IVES 2024) teve público recorde (246 participantes) e aprofundou o debate ao \textbf{propor} o tema “Beyond Borders: Evolving COIL Virtual Exchange” (Além das Fronteiras: Intercâmbio Virtual COIL em evolução).
O \textbf{foco} mudou da apresentação de experiências para a análise de \textbf{estratégias} de implementação, manutenção e \textbf{expansão} dos projetos, com destaque para novas abordagens pedagógicas e formatos inovadores de \textbf{interação}.
Os palestrantes internacionais \textbf{foram} Mirjam Hauck (Open University, Reino Unido); Andrea Alsina e Sofia Pekarek (Universidad Católica de Salta, Argentina); Sandra Valencia, (Universidad Católica de Manizales, Colômbia); Brenda García (Red LatAM COIL, México) e Gabriela Méndez (Florida International University, EUA).
A quarta edição (IVES 2025) reafirmou o evento como fórum internacional de \textbf{reflexão} estratégica, com a temática Changing Horizons (Mudança de Horizontes).
As discussões enfatizaram o \textbf{papel} dos \textbf{intercâmbios} virtuais frente às \textbf{transformações} educacionais \textbf{contemporâneas}, reunindo especialistas de diferentes países: João Monteiro, do ISP Gaya (Portugal), Stephanie Doscher, da University of Minnesota (EUA), Julia Bacca e Diana Garzon, da Uniminuto (Colômbia) e Adam Freed, da University of Michigan (EUA).
O evento evidenciou a \textbf{importância} da \textbf{troca} cultural e da formação global dos estudantes, além de dialogar com agendas mais amplas de internacionalização.
Ao \textbf{longo} dos anos, o IVES evoluiu de um \textbf{encontro} voltado à apresentação \textbf{institucional} para um espaço \textbf{consolidado} de debate global sobre Intercâmbios Virtuais.
Observa-se uma maior complexidade das discussões, passando da divulgação de práticas ao debate sobre \textbf{perspectivas} estratégicas e inovação pedagógica.
Essa trajetória reflete a \textbf{expansão} dos PCIs e o fortalecimento do Centro Paula Souza como \textbf{referência} internacional em COIL, iniciativa \textbf{essencial} para a internacionalização inclusiva e \textbf{colaborativa} no ensino \textbf{superior}.

% matched lemmas: ambiente, ampliar, anterior, apoio, aprender, basear, cada, coisa, colaborativo, compartilhar, competência, conectar, conexão, consolidar, construção, contemporâneo, contexto, coordenação, criar, crítico, desenvolvimento, destacar, diferença, disciplina, encontro, essencial, estratégia, estrutura, expandir, expansão, expectativa, foco, forma, fortalecer, futuro, gerar, habilidade, horário, humano, implementar, importância, iniciar, institucional, interação, intercâmbio, local, longo, mencionar, número, observar, papel, parte, pedagogia, permitir, perspectiva, pessoa, ponte, ponto, produzir, profissional, programa, promover, propor, prático, publicar, que, questão, quinto, reconhecer, reduzir, referência, reflexão, seguir, ser, significativo, superior, sustentável, tarefa, terceiro, todo, transformação, troca, verdadeiro
\end{textsample}
